\documentclass{ximera}

\title{Exam 2 Review}
\author{Rachel Bailey}
\license{CC: 0}         % replace with an appropriate license, or set it in xmPreamble

\begin{document}
\begin{abstract}
   This review covers covers sections...
\end{abstract}
\maketitle




\begin{problem}
Below are the graphs of $f(x), f'(x)$ and $f''(x)$. Identify each curve and explain your choices.
\end{problem}

\begin{problem}
    Use the \textbf{limit definition of the derivative} to find $f'(1)$ for $\displaystyle f(x)=\frac{1}{(x+2)^2}$.
    \vfill
\end{problem}

\begin{problem}
    Use the limit definition of the derivative to find $f'(x)$ for $\displaystyle f(x)=\frac{2}{\sqrt{2x-1}}$.
    \vfill
\end{problem}

\begin{problem}
    If the tangent line to $f(x)$ at $(4,3)$ passes through the point $(0,2)$, find $f(4)$ and $f'(4)$.
    \[
   f(4)= \answer{3}
    \]
    \[
    f'(4)=\answer{1/4}
    \]
\end{problem}

\begin{problem}
    Suppose the cost (in dollars) to produce $x$ units of a commodity is given by $C(x)$ where $C(10)=1000$ and $C'(10)=50$. Interpret the meanings of $C(10)$ and $C'(10)$ and find an estimate for the cost to produce 11 items.
    \[
    C(11)\approx \answer{1050}\text{ dollars}
    \]
\end{problem}

\begin{problem}
    Compute the the derivatives of the following functions:
\begin{enumerate}
    \item $\displaystyle F(x)=12x^2+\frac{4}{x^7}-\sqrt[3]{x^5}+5$
    \[
    F'(x)=\answer{24x-\frac{28}{x^8}-\frac{5}{3}x^{\frac{2}{3}}}
    \]

\item \(\displaystyle F(x)=2\tan(4^x)+\cot(3^x)\)
\[
F'(x)=\answer{2\sec^2(4^x)4^x\ln(4)-\csc^2(3^x)3^x\ln(3)}
\]

\item  $\displaystyle F(x)=\left(\frac{3}{x}+e^{x}-\sqrt{x^5}\right)\left(2x^{\pi}-\frac{x^2}{4}\right)$
\[
F'(x)=\answer{(-\frac{3}{x^2}+e^x-\frac{5}{2}x^{\frac{3}{2}})(2x^{\pi}-\frac{1}{4}x^2)+(2\pi x^{\pi-1}-\frac{1}{2}x)(\frac{3}{x}+e^x-\sqrt{x^5})}
\]

\item $\displaystyle F(x)=\frac{x^2+4\ln(x)-\sqrt[3]{x^4}}{x-1}$
\[
F'(x)=\answer{\frac{(2x+\frac{4}{x}-\frac{4}{3}x^{1/3}(x-1)-(x^2+4\ln(x)-x^{4/3})}{(x-1)^2}}
\]

\item $\displaystyle F(x)=10x\ln(2x^5-x)$
\[
F'(x)=\answer{10\ln(2x^5-x)+\frac{10x(10x^4-1)}{2x^5-x}}
\]

\item \( \displaystyle F(x)=\frac{4x-7}{x^2\sin(e^x)}\)
\[
F'(x)=\answer{\frac{4x^2\sin(e^x)-(2x\sin(e^x)+x^2e^x\cos(e^x))(4x-7)}{(x^2\sin(e^x))^2}}
\]

\item $\displaystyle F(x)=\frac{e^{x-\cos(x)}(x^7+6)}{(5-10x^7)^{10}}$
\[
F'(x)=\answer{\frac{(e^{x-\cos(x)}(1+\sin(x))(x^7+6)+7x^6e^{x-\cos(x)})(5-10x^7)^{10}+10(5-10x^7)^9(70x^6)(e^{x-\cos(x)}(x^7+6)}{(5-10x^7)^{20}}}
\]

\end{enumerate}
\end{problem}

\begin{problem}
Find the equation of the tangent line to \[f(x)=\frac{e^{2x^2+x}}{3x+1}\]at $x=0$.
\[
\text{Slope of tangent line: }m =\answer{-2}
\]
\[
\text{Point on the curve where the tangent line touches the curve: } (x,y)=(\answer{0},\answer{1})
\]
\[
\text{Equation of line in $y=mx+b$ form: } y=\answer{-2x+1}
\]
\end{problem}

\begin{problem}
Let $\displaystyle y=\frac{(x-1)\cos^2(x)}{x^2+1}$.
\begin{enumerate}
    \item Use \textbf{logarithmic differentiation} to find $\frac{dy}{dx}$.
    \[
    \frac{dy}{dx}=\answer{\frac{(x-1)\cos^2(x)}{x^2+1}(\frac{1}{x-1}-2\tan(x)-\frac{2x}{x^2+1})}
    \]
 
    \item Find the equation of the tangent line to $y$ at $x=0$. 
   
    \[
    \text{Slope at $x=0$: }m\answer{1}
    \]
    \[
    \text{Point on the curve where the tangent line touches the curve: }(x,y)=(\answer{0},\answer{-1})
    \]
    \[
    \text{Equation of the line in $y=mx+b$ form: }y=\answer{x-1}
    \]
\end{enumerate}
\end{problem}

\begin{problem}
    Use implicit differentiation to find $\frac{dy}{dx}$.
\begin{enumerate}

\item \(4x^3y+7y^2=10-12\sin(y)\)
\[
\frac{dy}{dx}=\answer{\frac{-12x^2y}{4x^3+14y+12\cos(y)}}
\]
\item \(4xe^{y}-ye^{2x}=4\)
\[
\frac{dy}{dx}=\answer{\frac{2e^{2x}y-4e^y}{4xe^y-e^{2x}}}
\]
\end{enumerate}
\end{problem}


\begin{problem}
    
Let $g(x)=xf(x)$. 
\begin{enumerate}
    \item What rule needs to be applied to find $g'(x)$?
    
   Rule: \wordChoice{
                \choice{Power Rule}
                \choice[correct]{Product Rule}
                \choice{Quotient Rule}
                \choice{Chain Rule} 
                }
    
    \item If $f(2)=4$ and $f'(2)=-1$, find $g'(2)$.
    \[
g'(2)=\answer{2}
\]
\end{enumerate}


\end{problem}

\begin{problem}
Compute $\displaystyle\frac{dy}{dx}$.
\begin{enumerate}
    \item $\displaystyle y= \cos^{-1}(x)\sin(x)$
\[
\frac{dy}{dx}=\answer{\frac{-\sin(x)}{\sqrt{1-x^2}+\cos(x)\arccos{x}}}
\]
\item $\displaystyle y=\tan^{-1}(e^x)$
\[
\frac{dy}{dx}=\answer{\frac{e^x}{1+e^{2x}}}
\]
\item $y=\displaystyle \frac{x}{\sin^{-1}(x)}$
\[
\frac{dy}{dx}=\answer{\frac{\arcsin{x}-\frac{x}{\sqrt{1-x^2}}}{(\arcsin{x})^2}}
\]
\end{enumerate}
\end{problem}

\begin{problem}
    Suppose a law-abiding chicken is crossing the road at a crosswalk a let its position (in inches) on the crosswalk be given by $s(t)=t^3-12t^2+21t$ where $t$ is in seconds ($t=0$ means the chicken is standing at the start of the crosswalk).
\begin{enumerate}
    \item At what time(s) $t$ is the chicken NOT moving?
    \[
    t=\answer{1} \text{ and }t=\answer{7}\text{seconds}
    \]
\item Is the chicken moving towards the other side of the road at $t=3$ or moving away?
\wordChoice{
                \choice{Towards}
                \choice[correct]{away}
            
                }
\item What is the chicken's acceleration at 5 seconds?
\[
a(4)=\answer{6}\text{in/sec}^2
\]

\end{enumerate}
\end{problem}

\begin{problem}
    Suppose is $V(t)$ the amount of air in a balloon at time $t$ (in seconds) and is given by
\[
V(t)=\frac{6\sqrt[3]{t}}{4t+1}.
\]Determine if the balloon is being filled with air, or drained of air at time $t=8$ seconds. \\


The volume is \wordChoice{
                \choice{increasing}
                \choice[correct]{decreasing}}
hence the balloon is being \wordChoice{
                \choice[correct]{drained}
                \choice{filled} }    of air.            
            

\end{problem}

\begin{problem}
    Let \(\displaystyle f(x)=(x-2)e^{\frac{x^2}{2}}\).  Find the $x$-values where the graph of $f(x)$ has a horizontal tangent.
    \[
    x=\answer{1}
    \]
\end{problem}

\begin{problem}
 Let $\displaystyle F(x)=f(x)g(x)$, $G(x)=\displaystyle f(g(x))$ and $H(x)=(x^2+1){g(x)}$. Use the following table to find $(a)-(d)$.
\begin{figure}[h!]
\centering
    \begin{tabular}{c|c|c|c|c|c|c}
        $x$&$f(x)$&$g(x)$& $f'(x)$&$g'(x)$&$f''(x)$&$g''(x)$ \\
        \hline
         2 &1&3& 2&-1&5&-4 \\
        \hline
         3 &-1&4& 10&0&3&-2
         \end{tabular}  

\end{figure}
\begin{enumerate}
    \item $F'(2)$.
    \[
    F'(2)=\answer{5}
    \]
\item $G'(2)$.
\[
G'(2)=\answer{-10}
\]
\item $H'(2)$.
\[
H'(2)=\answer{7}
\]
\item $F''(2)$.
\[
F''(2)=\answer{7}
\]

\end{enumerate}
\end{problem}

\begin{problem}
graph
\end{problem}

\begin{problem}
Use the linear approximation of $y=\sqrt[3]{x-1}$ at $x=2$ to approximate $\sqrt[3]{1.5}$ .
\[
\sqrt[3]{1.5}\approx \answer{7/6}
\]
\end{problem}

\begin{problem}
\begin{enumerate}
    \item Let $y=2x-4$. Compare the values of $\Delta y$ and $dy$ if $x$ changes from $1$ to $1.02$. 
    \[
    \Delta y=\answer{0.04}
    \]
    \[
    dy=\answer{0.04}
    \]

\item How can you explain your answer from part (a)?

\wordChoice{
                \choice{$\Delta y$ and $dy$ are always the same.}
                \choice{We are looking at a very small $dx$}
                \choice[correct]{The graph is a line, so it is equal to its tangent line at every $x$. } }
                

\end{enumerate}
\end{problem}
     

\end{document}